\documentclass[12pt,a4paper]{article}
\usepackage[utf8]{inputenc}
\usepackage[T2A]{fontenc}
\usepackage[russian]{babel}
\usepackage{amsmath,amssymb,amsthm}
\usepackage{booktabs,array}
\usepackage{microtype}
\usepackage{geometry}
\geometry{margin=2.3cm}
\setlength{\emergencystretch}{2em}

\newtheorem{proposition}{Предложение}
\newtheorem{nogocriterion}{Критерий провала}

\title{Аффинное меню spin-состояний и поколенческий триплет}
\author{}
\date{4 августа 2026 г.}

\begin{document}
\maketitle

\section{Смена принципа}

Геометрические автоморфизмы \(K=\mathbb{RP}^3\times S^1\) не делают три nonreference spin-сектора эквивалентными. Но исходная сумасбродная гипотеза утверждает, что \(K\) следует читать не буквально как пространство, а как меню состояний. Тогда естественной симметрией является не только группа диффеоморфизмов носителя, а полная группа аффинных relabelings spin-torsor.

Множество spin-структур является аффинной плоскостью
\[
\mathcal S_{spin}\cong\mathbb F_2^2
\]
из четырёх точек. Её полная affine group равна
\[
AGL(2,2)=\mathbb F_2^2\rtimes GL(2,2),
\qquad
|AGL(2,2)|=4\cdot6=24.
\]
Она действует верно на четырёх точках и изоморфна \(S_4\).

\section{Каноническое разложение \(1+3\)}

Рассмотрим комплексное пространство амплитуд четырёх menu-состояний:
\[
\mathcal H_{menu}=\mathbb C^4.
\]
Permutation representation \(S_4\) имеет каноническое разложение
\[
\mathbb C^4
=\mathbb C\mathbf 1
\oplus
\left\{v\in\mathbb C^4:\sum_{i=1}^4v_i=0\right\}.
\]
Первое слагаемое имеет размерность один, второе --- три. Соответствующие проекторы равны
\[
P_1=\frac14J,
\qquad
P_3=I-\frac14J,
\]
где \(J\) --- матрица из единиц.

\begin{proposition}[Канонический поколенческий carrier]
Полная affine symmetry четырёх spin-menu состояний единственным образом выделяет uniform singlet и трёхмерный sum-zero sector. Выбор одной конкретной spin-структуры не требуется.
\end{proposition}

\section{Уникальность invariant transition algebra}

Коммутант permutation representation имеет размерность два:
\[
\operatorname{End}_{S_4}(\mathbb C^4)
=\operatorname{span}\{I,J\}.
\]
Поэтому всякий fully affine-invariant Hermitian operator имеет вид
\[
T=aI+bJ
\]
и сохраняет то же разложение \(1+3\).

Два особенно простых оператора имеют ясный смысл. Laplacian полного графа
\[
L=4I-J
\]
имеет spectrum
\[
\operatorname{Spec}L=\{0,4,4,4\}.
\]
Uniform mode является нулевой, а triplet --- вырожденным excited sector. Наоборот, rank-one operator
\[
M=J
\]
имеет spectrum
\[
\operatorname{Spec}M=\{0,0,0,4\}
\]
и делает тяжёлым только uniform singlet, оставляя canonical rank-three kernel.

\section{Соединение с \(SU(5)\)}

Предыдущий gate построил одно anomaly-free поколение как
\[
\mathbf{10}\oplus\overline{\mathbf5}.
\]
Теперь внутреннее matter space можно записать как
\[
\mathcal H_{matter}
=\operatorname{im}P_3
\otimes
\left(\mathbf{10}\oplus\overline{\mathbf5}\right).
\]
Его размерность равна
\[
3\cdot(10+5)=45,
\]
то есть оно содержит три одинаковых chiral поколения. Поскольку anomaly одного package равна нулю, anomaly трёх копий также равна нулю.

\begin{proposition}[Три поколения без выбора трёх точек]
Если полная affine group spin-menu является физической symmetry переходов, то tensor product её canonical triplet с одним \(SU(5)\) chiral package даёт три эквивалентных anomaly-free поколения.
\end{proposition}

\section{Что здесь действительно новое}

Механизм не использует формулу \(4-1=3\) и не выбирает одну spin-структуру как лишнюю. Поколения являются не тремя отдельными точками меню, а тремя независимыми sum-zero superpositions четырёх spin-состояний. Uniform superposition играет роль singlet/reference mode.

Это устраняет прежнюю проблему affine torsor: reference определяется как invariant vector, а не как выбранный origin.

\section{Новый no-go}

Полная affine symmetry не следует из диффеоморфизмов исходного \(K\). Она является новым принципом именно menu-интерпретации. Кроме того, точная \(S_4\)-симметрия делает три поколения вырожденными и не объясняет mass hierarchy или mixing.

\begin{nogocriterion}[Провал завершённой family theory]
Разложение \(1+3\) выводит число и эквивалентность поколений, но не их разные массы. Физическая версия должна вывести symmetry-breaking operator из уже существующего различия между torsion- и circle-направлениями, не вводя произвольную Yukawa matrix.
\end{nogocriterion}

\section{Следующий gate}

Геометрическая subgroup, сохраняющая integral decomposition \(\mathbb Z_2\times\mathbb Z\), уже разбивает три labels по схеме \(1+2\). Следующий тест должен спроецировать этот единственный доступный breaking operator на canonical triplet \(\operatorname{im}P_3\) и вычислить его spectrum.

Если получится фиксированный singlet--doublet family pattern, menu-модель впервые даст не только число поколений, но и форму первого mass splitting. Если же projected operator содержит свободные коэффициенты или не различает family modes, ветвь закроется на уровне красивой групповой конструкции.

\end{document}